\documentclass[conference]{IEEEtran}
\IEEEoverridecommandlockouts
% The preceding line is only needed to identify funding in the first footnote. If that is unneeded, please comment it out.
\usepackage{cite}
\usepackage{amsmath,amssymb,amsfonts}
\usepackage{algorithmic}
\usepackage{graphicx}
\usepackage{textcomp}
\usepackage{xcolor}
\def\BibTeX{{\rm B\kern-.05em{\sc i\kern-.025em b}\kern-.08em
    T\kern-.1667em\lower.7ex\hbox{E}\kern-.125emX}}  
\begin{document}

\title{Project Paper Title*\\
{\footnotesize \textsuperscript{*}Note: Sub-titles are not captured in Xplore and
should not be used}
\thanks{Identify applicable funding agency here. If none, delete this.}
}

\author{\IEEEauthorblockN{1\textsuperscript{st} Given Name Surname}
\IEEEauthorblockA{\textit{dept. name of organization (of Aff.)} \\
\textit{name of organization (of Aff.)}\\
City, Country \\
email address or ORCID}
\and
\IEEEauthorblockN{2\textsuperscript{nd} Given Name Surname}
\IEEEauthorblockA{\textit{dept. name of organization (of Aff.)} \\
\textit{name of organization (of Aff.)}\\
City, Country \\
email address or ORCID}
\and
\IEEEauthorblockN{3\textsuperscript{rd} Given Name Surname}
\IEEEauthorblockA{\textit{dept. name of organization (of Aff.)} \\
\textit{name of organization (of Aff.)}\\
City, Country \\
email address or ORCID}
\and
\IEEEauthorblockN{4\textsuperscript{th} Given Name Surname}
\IEEEauthorblockA{\textit{dept. name of organization (of Aff.)} \\
\textit{name of organization (of Aff.)}\\
City, Country \\
email address or ORCID}
\and
\IEEEauthorblockN{5\textsuperscript{th} Given Name Surname}
\IEEEauthorblockA{\textit{dept. name of organization (of Aff.)} \\
\textit{name of organization (of Aff.)}\\
City, Country \\
email address or ORCID}
\and
\IEEEauthorblockN{6\textsuperscript{th} Given Name Surname}
\IEEEauthorblockA{\textit{dept. name of organization (of Aff.)} \\
\textit{name of organization (of Aff.)}\\
City, Country \\
email address or ORCID}
}

\maketitle

\begin{abstract}

\noindent[Write 1/2 sentences about the project problem you are solving.]

\noindent[Write 2/3 sentences about the solution that you have implemented.]

\noindent[Write 1/2 sentences about the implemented system's performance/accuracy/efficiency.]
\end{abstract}

\begin{IEEEkeywords}
component, formatting, style, styling, insert
\end{IEEEkeywords}

\section{Introduction}

\noindent Please write this section according to the following instructions:

\noindent \textbf{[Problem Statement]} In this part, you need to explain the idea of your project in detail. Also, clearly describe the motivations behind your project's problem and explain why you chose to implement this problem.

\noindent \textbf{[Literature Review]} Search google, google scholar or research gate for the problems/papers related to your projects. Discuss the previous works done related to your projects and what are the features implemented in these previous works.

\noindent \textbf{[Innovative features added to your projects]} Describe how your project differs from existing ideas. Write the summary of your project's features. Add what features you have added which did not exist in the previous similiar type of projects.

\section{Proposed Method}

\subsection{Block diagram of the overall system}

You have implemented a hardware system in your project and the system is functioning with the software codes you have uploaded to your arduino or raspberry pi. Now, in this section you have to present a block diagram of the overall system of your project. You should prepare the block diagram using microsoft powerpoint or other tools. The block diagram must show each sub part such as sensors, motors, arduino etc. of the whole system and how they are connected with each other.

\noindent For example: let's say you have made a project where a gas sensor detects gas in a kitchen and if the gas level goes beyond a limit, an alaram will ring automatically by the arduino. For such simple system, you must connect the gas sensor and the alarm to the microcontroller. Gas sensor reads the concentration level of gas in the kitchen and the microcontroller continously reads this concentration. When the concentration goes beyond the limit (you have set the limit in the code of the microcontroller), the microcontroller turns on the alaram. Now, you can understand the block diagram of your system easily. There will be three block diagrams: microcontroller, gas sensor and alarm module. Gas sensor block diagram will be connected to microcontroller and Alaram module's block diagram will also be connected to the microcontroller. You should put additional texts on the block diagram in a way so that the readers (even non technical people) can understand how your system is operating just by looking at the block diagram.

\noindent Now, you also have to explain in details about your block diagram. You must write about each block diagram of the enrite system in details and how they are connected to each other and how they are working.



\section{Implemented Hardware System }
\subsection{List of the hardwares and software Environment}

In this section, you have to describe the hardware equipment used in this project. Each hardware equipment and their model number must be mentioned.

\noindent Also, you have to mention the software environment you have used.

\subsection{Figures of the implemented hardware system}\label{AA}
In this section, add the actual image of your project. Don't just put an image you captured, rather indicate on the figure the main parts of the project. It may help if you can add multiple images so that anyone can understand the entire hardware system implemented.

\noindent Describe breifly about the figures you added in this subsection.


\subsection{Figures and Tables}

\noindent \textbf{If you want to add figures or tables, use the syntax of the codes from this section!}

\paragraph{Positioning Figures and Tables} Place figures and tables at the top and 
bottom of columns. Avoid placing them in the middle of columns. Large 
figures and tables may span across both columns. Figure captions should be 
below the figures; table heads should appear above the tables. Insert 
figures and tables after they are cited in the text. Use the abbreviation 
``Fig.~\ref{fig}'', even at the beginning of a sentence.

\begin{table}[htbp]
\caption{Table Type Styles}
\begin{center}
\begin{tabular}{|c|c|c|c|}
\hline
\textbf{Table}&\multicolumn{3}{|c|}{\textbf{Table Column Head}} \\
\cline{2-4} 
\textbf{Head} & \textbf{\textit{Table column subhead}}& \textbf{\textit{Subhead}}& \textbf{\textit{Subhead}} \\
\hline
copy& More table copy$^{\mathrm{a}}$& &  \\
\hline
\multicolumn{4}{l}{$^{\mathrm{a}}$Sample of a Table footnote.}
\end{tabular}
\label{tab1}
\end{center}
\end{table}

\begin{figure}[htbp]
\centerline{\includegraphics{fig1.png}}
\caption{Example of a figure caption.}
\label{fig}
\end{figure}

Figure Labels: Use 8 point Times New Roman for Figure labels. Use words 
rather than symbols or abbreviations when writing Figure axis labels to 
avoid confusing the reader. As an example, write the quantity 
``Magnetization'', or ``Magnetization, M'', not just ``M''. If including 
units in the label, present them within parentheses. Do not label axes only 
with units. In the example, write ``Magnetization (A/m)'' or ``Magnetization 
\{A[m(1)]\}'', not just ``A/m''. Do not label axes with a ratio of 
quantities and units. For example, write ``Temperature (K)'', not 
``Temperature/K''.

\section{Results}

In this section, describe the output results of your system. If we take the example of the gas sensor and fire alarm, then as output results, you should add:

\noindent 1) Gas sensor's data reading figure when there is no gas leakage and Gas sensor's data reading figure when there is certain amount of gas leakage.

\noindent 2) When the concentration of the gas is beyond a limit, the microcontroller must send a signal, show in a figure/code, how that is happening.

\noindent 3) Finally, show that the alarm is turning ON in a figure.

\noindent You may additional info such as respond time of the alarm to turn ON after detecting gas leakage and other important attributes of your overall system.

\noindent We described this results section according to a simple project mentioned above. But for your project, you may have different sensors and setups. So, you must carefully add step by step results so that we can understand each part's output and assess properly how good your system is working.
\section{Discussion}
In this section, you should mention about the limitations of your project. Show one or more example cases where your project will not work as expected.

While implementing the solution, there must be various requirements in your projects. While implementing these requirements, did you face any restrictions? Are the restrictions resulting into trade-off decision making (performance, cost, etc)? For example, let's say you want to design a hardware system that produces decision based output very quickly by processing a large amount of data. To do that, you either can use a high performance microprocessors which is costly or you may use other services or algorithm that can process the large amount data efficiently while using the low performance microprocessors. The second approach (better algorithm/technique in low performance microprocessors/devices) may not be faster than the first approach (high performance microprocessors/devices) but the second approach is cost-effective. If cost-effective is a requirement in your projects, then it conflicts the first approach. So, you need to consider the trade-off here while implementing your projects in such scenarios. Write 2/3 sentences if you face such conflicting requirements.  (\textbf{P2}) 

Think of two alternate solution/design other than the one you have implemented. Describe why or how they are better or worse or different (budget, performance, etc) than your current implementation (\textbf{P3}).

What will be the impact of this solution? Did you solve a problem from Computer Science field that an CSE student faces regularly. Or your problem has got something to do with other fields like agriculture, medicine, environment, education, etc something that you dont encounter regularly but now you have to know/interact with something outside your  familiar area of knowledge (\textbf{P4}).

Have you studied any communication standard, interfacing standard documents, manuals when implementing the proposed solution? Did everything went according to the manual or the standard documents? Or did you have any experience outside the standards? Also, did you read the ethical sides of your projects features while implementing them and maintained the ethical code of conducts? Write 2-3 sentences on this if applicable (\textbf{P5}).

Have you talked to any stakeholders (users/managers/funders/etc) who might be interested in your project? If yes, then write how their requirements were incorporated in your solution or design (\textbf{P6}).

How decoupled are the modules that you have in your system? How dependent are they? Is there hardware and software modules dependent on each other and yet not on the same device? Write 2-3 sentences on this aspect (\textbf{P7}). 



\begin{table}[h]
    \begin{center}
    \begin{tabular}{|c|c|c|c|c|c|c|}
\hline  
P1&P2&P3&P4&P5&P6&P7\\
\hline 
$\checkmark$&&&&&&\\
\hline 
\end{tabular}
\end{center}
    \caption{Complex Engineering Problem Mapping. Tic P1 and some or more from P2-P7 as applicable. P1 - Depth of Knowledge, P2 - Conflicting Requirements, P3 - Depth of Analysis, P4 - Familiarity of Issues, P5 - Extent of Applicable Codes, P6 - Extent of Stakeholders, P7 - Interdependence}
    \label{tab:cep}
\end{table}

\section{Conclusion}

Finally, in this section, discuss briefly about your project's features in 1/2 sentences. Then, discuss your system's efficiency and what are the future scopes of your project in 3/4 sentences!

\section*{Acknowledgment}

\textbf{If you want to mention who helped you throughout this project you should mention here. Also, if you used the equipments from the lab or a friend, you should mention that also.}


\noindent The preferred spelling of the word ``acknowledgment'' in America is without 
an ``e'' after the ``g''. Avoid the stilted expression ``one of us (R. B. 
G.) thanks $\ldots$''. Instead, try ``R. B. G. thanks$\ldots$''. Put sponsor 
acknowledgments in the unnumbered footnote on the first page.

\section*{References}

\textbf{[You must add the github link and youtube video link in the reference link]}

\noindent Please number citations consecutively within brackets \cite{b1}. The 
sentence punctuation follows the bracket \cite{b2}. Refer simply to the reference 
number, as in \cite{b3}---do not use ``Ref. \cite{b3}'' or ``reference \cite{b3}'' except at 
the beginning of a sentence: ``Reference \cite{b3} was the first $\ldots$''

Number footnotes separately in superscripts. Place the actual footnote at 
the bottom of the column in which it was cited. Do not put footnotes in the 
abstract or reference list. Use letters for table footnotes.

Unless there are six authors or more give all authors' names; do not use 
``et al.''. Papers that have not been published, even if they have been 
submitted for publication, should be cited as ``unpublished'' \cite{b4}. Papers 
that have been accepted for publication should be cited as ``in press'' \cite{b5}. 
Capitalize only the first word in a paper title, except for proper nouns and 
element symbols.

For papers published in translation journals, please give the English 
citation first, followed by the original foreign-language citation \cite{b6}.

\begin{thebibliography}{00}
\bibitem{b1} G. Eason, B. Noble, and I. N. Sneddon, ``On certain integrals of Lipschitz-Hankel type involving products of Bessel functions,'' Phil. Trans. Roy. Soc. London, vol. A247, pp. 529--551, April 1955.
\bibitem{b2} J. Clerk Maxwell, A Treatise on Electricity and Magnetism, 3rd ed., vol. 2. Oxford: Clarendon, 1892, pp.68--73.
\bibitem{b3} I. S. Jacobs and C. P. Bean, ``Fine particles, thin films and exchange anisotropy,'' in Magnetism, vol. III, G. T. Rado and H. Suhl, Eds. New York: Academic, 1963, pp. 271--350.
\bibitem{b4} K. Elissa, ``Title of paper if known,'' unpublished.
\bibitem{b5} R. Nicole, ``Title of paper with only first word capitalized,'' J. Name Stand. Abbrev., in press.
\bibitem{b6} Y. Yorozu, M. Hirano, K. Oka, and Y. Tagawa, ``Electron spectroscopy studies on magneto-optical media and plastic substrate interface,'' IEEE Transl. J. Magn. Japan, vol. 2, pp. 740--741, August 1987 [Digests 9th Annual Conf. Magnetics Japan, p. 301, 1982].
\bibitem{b7} M. Young, The Technical Writer's Handbook. Mill Valley, CA: University Science, 1989.
\end{thebibliography}
\vspace{12pt}
\color{red}
\end{document}
